\documentclass{beamer}
\usepackage[utf8]{inputenc}
\usepackage[T1]{fontenc}
\usepackage[portuguese]{babel}
\usepackage{amsmath}
\usepackage{amssymb}
\usepackage{tikz}
\usetikzlibrary{shapes.geometric, arrows.meta, calc, decorations.pathreplacing, patterns}

% Configuração do Tema
\usetheme{Madrid}
\usecolortheme{default}
\beamertemplatenavigationsymbolsempty

% Comando para o círculo vermelho
\providecommand{\circled}[1]{\mbox{\tikz[baseline=(char.base)]{\node[shape=circle,draw=red,thick,inner sep=0.5pt,text=red,font=\tiny\bfseries] (char) {#1};}}}

% Convenções visuais
\tikzset{
	eixo/.style={->, thick},
	vetorf/.style={->, thick, orange!90!black},
	vetord/.style={->, thick, blue!75!black},
	vetorv/.style={->, thick, magenta!80!black},
	vetorg/.style={->, thick, red!80!black},
	vetors/.style={->, thick, green!60!black},
	ponto/.style={circle, fill=black, inner sep=1.5pt}
}

% Informações do Título
\title{Aula: Teoria e Prática}
\subtitle{Sistemas de Partículas: Centro de Massa, Momento Linear e Colisões}
\date{}

\begin{document}
	
	% SLIDE 1
	\frame{\titlepage}
	
	
	% SLIDE 2
	\begin{frame}{De uma partícula a um sistema de partículas}
		\scriptsize
		
		Até aqui descrevemos muitas situações tratando o corpo como uma partícula. Agora consideramos um sistema com
		$
		n>1
		$
		partículas, com massas \(m_1,m_2,\ldots,m_n\) e posições
		$
		\vec r_1,\vec r_2,\ldots,\vec r_n .
		$
		Para cada partícula, a segunda lei de Newton pode ser escrita como
		\[
		\textcolor{red}{m_i\vec a_i=\vec F_i, \qquad i=1,2,\ldots,n.}
		\]
		
		\begin{columns}[T]
			\begin{column}{0.7\textwidth}
				\textcolor{blue}{\textbf{Problema:}}
				
				Para \(n>2\), resolver todas as equações de Newton pode ser difícil.
				
				\vspace{0.15cm}
				
				\textcolor{blue}{\textbf{Ideia principal:}}
				
				Muitas vezes não precisamos conhecer a trajetória de cada partícula. Basta estudar um ponto especial do sistema:
				\[
				\textcolor{red}{\text{centro de massa.}}
				\]
				
				Esse ponto resume o movimento global do sistema.
			\end{column}
			
			\begin{column}{0.3\textwidth}
				\centering
				\begin{tikzpicture}[scale=0.85,>=stealth]
					\draw[eixo] (-0.3,0) -- (3.6,0) node[right] {$x$};
					\draw[eixo] (0,-0.3) -- (0,2.7) node[above] {$y$};
					
					\filldraw[fill=blue!25, draw=blue!80!black] (0.8,1.9) circle (0.13);
					\node[above] at (0.8,2.05) {\scriptsize \(m_1\)};
					\filldraw[fill=green!25, draw=green!60!black] (2.7,1.7) circle (0.18);
					\node[above] at (2.7,1.9) {\scriptsize \(m_2\)};
					\filldraw[fill=orange!35, draw=orange!80!black] (1.8,0.7) circle (0.11);
					\node[below] at (1.8,0.55) {\scriptsize \(m_3\)};
					
					\filldraw[fill=red!70!black] (1.9,1.45) circle (0.06);
					\node[right] at (1.98,1.45) {\scriptsize CM};
					
					\draw[vetord] (0,0) -- (0.8,1.9) node[midway,left] {\scriptsize \(\vec r_1\)};
					\draw[vetord] (0,0) -- (2.7,1.7) node[midway,above] {\scriptsize \(\vec r_2\)};
					\draw[vetord] (0,0) -- (1.8,0.7) node[midway,below] {\scriptsize \(\vec r_3\)};
				\end{tikzpicture}
			\end{column}
		\end{columns}
		
	\end{frame}
	
	
	% SLIDE 3
	\begin{frame}{Centro de massa (CM): definição e exemplo em 1D}
		\tiny
		
		\textbf{Definição.} Para um sistema de partículas com massa total
		\(\textcolor{red}{M=m_1+m_2+\cdots+m_n}\), o centro de massa é
		\[
		\textcolor{red}{
			\vec r_{\text{cm}}=
			\frac{m_1\vec r_1+m_2\vec r_2+\cdots+m_n\vec r_n}{M}}.
		\]
		
		No plano:
		$
		\textcolor{red}{
			x_{\text{cm}}=\frac{m_1x_1+\cdots+m_nx_n}{M}},
		\qquad
		\textcolor{red}{
			y_{\text{cm}}=\frac{m_1y_1+\cdots+m_ny_n}{M}}.
		$
		
		\vspace{0.3cm}
		
		\begin{columns}[T]
			\begin{column}{0.50\textwidth}
				\textcolor{blue}{\textbf{Leitura física}}
				\begin{itemize}
					\item \(\vec r_{\text{cm}}\) é uma média ponderada pelas massas;
					\item massas maiores deslocam o CM para mais perto delas;
					\item o CM não precisa coincidir com uma partícula real.
				\end{itemize}
				
				\vspace{0.05cm}
				
				\textcolor{blue}{\textbf{Exemplo: duas partículas em 1D}}
				
				\vspace{0.3cm}
				
				Para \(x_1=0\) e \(x_2=d\):
				\[
				\textcolor{red}{
					x_{\text{cm}}=
					\frac{M_1x_1+M_2x_2}{M_1+M_2}
					=
					\frac{M_2d}{M_1+M_2}.}
				\]
				
				\vspace{-0.2cm}
				
				\[
				\begin{array}{c|c}
					\text{condição} & \text{posição do CM} \\ \hline
					M_1=M_2 & x_{\text{cm}}=d/2 \\[1pt]
					M_1\gg M_2 & x_{\text{cm}}\approx 0
				\end{array}
				\]
				
				\textcolor{red}{\textbf{Mensagem:} o CM fica mais perto da massa maior.}
			\end{column}
			
			\begin{column}{0.50\textwidth}
				\centering
				\begin{tikzpicture}[scale=0.82,>=stealth]
					\draw[eixo] (-0.4,0) -- (4.3,0) node[right] {$x$};
					
					% massas
					\filldraw[fill=blue!25, draw=blue!80!black] (0,0) circle (0.24);
					\node[below] at (0,-0.30) {\scriptsize \(M_1\)};
					\node[above] at (0,0.28) {\scriptsize \(x_1=0\)};
					
					\filldraw[fill=green!25, draw=green!60!black] (3.3,0) circle (0.14);
					\node[below] at (3.3,-0.30) {\scriptsize \(M_2\)};
					\node[above] at (3.3,0.28) {\scriptsize \(x_2=d\)};
					
					% distância
					\draw[decorate, decoration={brace, mirror, amplitude=5pt}] (0,-0.72) -- (3.3,-0.72);
					\node[below] at (1.65,-0.93) {\scriptsize \(d\)};
					
					% CM
					\filldraw[fill=red!80!black] (0.55,0) circle (0.065);
					\node[above, red!80!black] at (0.55,0.85) {\scriptsize CM};
					\draw[dashed, red!80!black] (0.55,0.08) -- (0.55,0.9);
					
					% caso limite
					\node[align=center, font=\scriptsize, text=red!80!black] at (2.25,1.05)
					{\(M_1\gg M_2\)\\ \(x_{\text{cm}}\approx 0\)};
				\end{tikzpicture}
				
				\vspace{0.15cm}
				
				\[
				\textcolor{red}{
					M_1=M_2 \Rightarrow x_{\text{cm}}=\frac d2,
					\qquad
					M_1\gg M_2 \Rightarrow x_{\text{cm}}\approx 0.}
				\]
			\end{column}
		\end{columns}
		
	\end{frame}
	
	
	% SLIDE 4
	\begin{frame}{Movimento do centro de massa}
		\scriptsize
		
		Partindo da definição do centro de massa e derivando em relação ao tempo:
		\[
		\textcolor{red}{
			\vec v_{\text{cm}}
			=
			\frac{d\vec r_{\text{cm}}}{dt}
			=
			\frac{m_1\vec v_1+m_2\vec v_2+\cdots+m_n\vec v_n}{M}.}
		\]
		
		Derivando novamente:
		\[
		\textcolor{red}{
			\vec a_{\text{cm}}
			=
			\frac{d\vec v_{\text{cm}}}{dt}
			=
			\frac{m_1\vec a_1+m_2\vec a_2+\cdots+m_n\vec a_n}{M}.}
		\]
		
		Multiplicando por \(M\):
		$
		{
			M\vec a_{\text{cm}}
			=
			m_1\vec a_1+m_2\vec a_2+\cdots+m_n\vec a_n.}
		$
		Pela segunda lei de Newton, \(m_i\vec a_i=\vec F_i\). Logo:
		\[
		\textcolor{red}{
			M\vec a_{\text{cm}}
			=
			\vec F_1+\vec F_2+\cdots+\vec F_n
			=
			\vec F_{\text{ext}}.}
		\]
		
		\textcolor{blue}{\textbf{Ideia física:}}
		
		As forças internas cancelam-se em pares. Portanto, o movimento global do centro de massa é determinado apenas pela força externa resultante:
		\[
		\textcolor{red}{M\vec a_{\text{cm}}=\vec F_{\text{ext}}.}
		\]
		
		Se
		$
		\vec F_{\text{ext}}=0,
		$
		então
		$
		\vec a_{\text{cm}}=0
		\qquad \Rightarrow \qquad
		\vec r_{\text{cm}}=\vec r_0+\vec v_{\text{cm}}t.
		$
		Ou seja: o centro de massa permanece em repouso ou move-se com velocidade constante.
		
	\end{frame}
	
	
	% SLIDE 5
	\begin{frame}{Momento linear de uma partícula e de um sistema}
		\tiny
		
		\textcolor{blue}{\textbf{Partícula.}}
		Para uma partícula de massa \(m\), com velocidade \(\vec v\), o momento linear é
		$
		\textcolor{red}{\vec p=m\vec v},
		\qquad
		[\vec p]=\textcolor{red}{\mathrm{kg\,m/s}}.
		$
		
		\begin{columns}[T]
			\begin{column}{0.66\textwidth}
				\textcolor{blue}{\textbf{Sistema de partículas.}}
				Se o sistema tem \(n\) partículas, o momento linear total é a soma vetorial:
				\[
				\textcolor{red}{
					\vec P=\vec p_1+\vec p_2+\cdots+\vec p_n
					=
					m_1\vec v_1+m_2\vec v_2+\cdots+m_n\vec v_n.}
				\]
				
				Como
				$
				\vec v_{\text{cm}}=
				\frac{m_1\vec v_1+m_2\vec v_2+\cdots+m_n\vec v_n}{M},
				$
				então
				$
				\textcolor{red}{\vec P=M\vec v_{\text{cm}}.}
				$
				
				\vspace{0.25cm}
				
				\textcolor{blue}{\textbf{Leitura física:}}
				o momento linear total do sistema é o momento linear de uma partícula equivalente,
				com massa \(M\), movendo-se com a velocidade do centro de massa.
			\end{column}
			
			\begin{column}{0.34\textwidth}
				\centering
				\begin{tikzpicture}[scale=0.62,>=stealth]
					\draw[fill=blue!20, draw=blue!80!black] (0,1.8) circle (0.13);
					\draw[vetorv] (0.15,1.8) -- (1.2,2.15)
					node[midway, above] {\scriptsize \(\vec p_1\)};
					
					\draw[fill=green!20, draw=green!60!black] (1.2,0.6) circle (0.16);
					\draw[vetorv] (1.35,0.6) -- (2.3,0.25)
					node[midway, below] {\scriptsize \(\vec p_2\)};
					
					\draw[fill=orange!30, draw=orange!80!black] (2.4,1.5) circle (0.11);
					\draw[vetorv] (2.55,1.5) -- (3.1,2.2)
					node[midway, right] {\scriptsize \(\vec p_3\)};
					
					\filldraw[fill=red!80!black] (1.45,1.35) circle (0.06);
					\draw[vetorf] (1.45,1.35) -- (2.75,1.6)
					node[midway, above] {\scriptsize \(\vec P\)};
					
					\node[align=center, font=\scriptsize] at (1.55,-0.65)
					{\tiny \(\vec P\) é a soma vetorial\\ \tiny dos momentos lineares};
				\end{tikzpicture}
			\end{column}
		\end{columns}
		
		\vspace{0.1cm}
		
		\begin{center}
			\fbox{
				\begin{minipage}{0.94\textwidth}
					\tiny
					\textcolor{blue}{\textbf{Segunda lei de Newton escrita com o momento linear.}}
					Para cada partícula:
					$
					\textcolor{red}{
						\vec F_i=\frac{d\vec p_i}{dt}}
					\qquad
					(i=1,2,\ldots,n).
					$
					
					Como \(\vec P=\vec p_1+\vec p_2+\cdots+\vec p_n\), temos
					\[
					\frac{d\vec P}{dt}
					=
					\frac{d\vec p_1}{dt}
					+
					\frac{d\vec p_2}{dt}
					+\cdots+
					\frac{d\vec p_n}{dt}
					=
					\vec F_1+\vec F_2+\cdots+\vec F_n.
					\]
					
					As forças internas do sistema cancelam-se em pares; assim,
					\(\vec F_1+\vec F_2+\cdots+\vec F_n=\vec F_{\text{ext}}\).
					Logo,
					\[
					\textcolor{red}{
						\frac{d\vec P}{dt}=\vec F_{\text{ext}}.}
					\]
					
					\vspace{-0.15cm}
					
					\textcolor{red}{\textbf{Mensagem:}}
					o momento linear total de um sistema só muda se existir força externa resultante.
			\end{minipage}}
		\end{center}
		
	\end{frame}
	
	
	% SLIDE 6
	\begin{frame}{Conservação do momento linear, colisão e impulso}
		\tiny
		
		\textcolor{blue}{\textbf{Ideia central.}}
		Para um sistema de partículas:
		$
		\textcolor{red}{
			\frac{d\vec P}{dt}=\vec F_{\text{ext}}.}
		$
		
		\vspace{0.2cm}
		
		\begin{columns}[T]
			\begin{column}{0.52\textwidth}
				\textcolor{blue}{\textbf{Conservação do momento linear}}
				
				\vspace{0.2cm}
				
				Se a força externa resultante é nula,
				$
				\vec F_{\text{ext}}=0,
				$
				então o momento linear total não muda:
				\[
				\textcolor{red}{\vec P=\text{constante}}
				\qquad \Rightarrow \qquad
				\textcolor{red}{\vec P_i=\vec P_f.}
				\]
				
				Em componentes:
				$
				\textcolor{red}{
					P_{x,i}=P_{x,f}, \qquad
					P_{y,i}=P_{y,f}.}
				$
				
				\vspace{0.2cm}
				
				\textcolor{blue}{\textbf{Uso em colisões.}}
				Durante uma colisão de curta duração, se a força externa resultante é desprezável,
				usamos \(\vec P_i=\vec P_f\) para relacionar as velocidades antes e depois.
			\end{column}
			
			\begin{column}{0.48\textwidth}
				\textcolor{blue}{\textbf{Colisão e impulso}}
				
				\vspace{0.2cm}
				
				Durante uma colisão, a velocidade pode mudar, em módulo e sentido, num intervalo curto \(\Delta t\).
				A variação do momento linear é
				\[
				\textcolor{red}{\Delta \vec p=\vec p_f-\vec p_i.}
				\]
				
				Essa variação é chamada impulso:
				$
				\textcolor{red}{\vec J=\Delta \vec p.}
				$
				
				Se \(F\) é o módulo da força média durante a colisão:
				\[
				\textcolor{red}{
					F=\frac{\Delta p}{\Delta t}},
				\qquad
				\Delta p=|\Delta \vec p|.
				\]
				
				\textcolor{red}{\textbf{Interpretação:}}
				para a mesma \(\Delta p\),
				\[
				\textcolor{red}{
					\Delta t \text{ maior}
					\quad \Rightarrow \quad
					F \text{ menor}.}
				\]
			\end{column}
		\end{columns}
		
		\vspace{0.05cm}
		
		\begin{columns}[T]
			\begin{column}{0.52\textwidth}
				\vspace{-1.15cm}
				
				\centering
				\begin{tikzpicture}[scale=0.82,>=stealth]
					% antes
					\node[font=\tiny] at (0,1.95) {antes};
					\draw[fill=blue!20, draw=blue!80!black] (-0.95,1.35) circle (0.16);
					\draw[fill=green!20, draw=green!60!black] (0.95,1.35) circle (0.16);
					\draw[vetorv] (-0.95,1.35) -- (-0.15,1.35) node[midway, above] {\tiny \(\vec v_{1i}\)};
					\draw[vetorv] (0.95,1.35) -- (0.25,1.35) node[midway, above] {\tiny \(\vec v_{2i}\)};
					
					% depois
					\node[font=\tiny] at (0,0.45) {depois};
					\draw[fill=blue!20, draw=blue!80!black] (-0.45,-0.10) circle (0.16);
					\draw[fill=green!20, draw=green!60!black] (0.45,-0.10) circle (0.16);
					\draw[vetorv] (-0.45,-0.10) -- (-1.20,-0.10) node[midway, below] {\tiny \(\vec v_{1f}\)};
					\draw[vetorv] (0.45,-0.10) -- (1.20,-0.10) node[midway, below] {\tiny \(\vec v_{2f}\)};
					
					\node[align=center, font=\tiny, red!80!black] at (0,-0.95)
					{\(\vec P_i=\vec P_f\) se \(\vec F_{\text{ext}}=0\)};
				\end{tikzpicture}
			\end{column}
			
			\begin{column}{0.48\textwidth}
				\centering
				\begin{tikzpicture}[scale=0.86,>=stealth]
					% chão
					\draw[thick] (-1.8,-0.45) -- (2.0,-0.45);
					
					% parede/alvo
					\draw[fill=gray!20, draw=black, pattern=north east lines] (0.95,-0.45) rectangle (1.25,1.05);
					\node[right] at (1.25,0.25) {\tiny alvo};
					
					% partícula antes
					\draw[fill=blue!25, draw=blue!80!black] (-1.05,0.00) circle (0.14);
					\draw[vetorv] (-1.05,0.00) -- (-0.20,0.00) node[midway, above] {\tiny \(\vec p_i\)};
					
					% partícula no instante logo após o choque
					\draw[fill=blue!25, draw=blue!80!black] (0.75,0.00) circle (0.14);
					\draw[vetorv] (0.75,0.00) -- (-0.00,0.00) node[midway, below] {\tiny \(\vec p_f\)};
					
					% delta p junto da parede
					\draw[vetorf] (0.90,0.42) -- (0.15,0.42) node[midway, above] {\tiny \(\Delta \vec p\)};
					\draw[dashed] (0.90,0.42) -- (0.90,0.05);
					
					\node[align=center, font=\tiny, red!80!black] at (0.1,-1.05)
					{choque no alvo: \\ \(\Delta \vec p\) ocorre no contacto com a parede};
				\end{tikzpicture}
			\end{column}
		\end{columns}
		
	\end{frame}
	
	
	% SLIDE 7
	\begin{frame}{Como resolver problemas de colisões}
		\scriptsize
		
		\begin{columns}[T]
			\begin{column}{0.57\textwidth}
				\textcolor{blue}{\textbf{Procedimento recomendado:}}
				
				\begin{enumerate}
					\item Escolher o sistema: quais corpos fazem parte da colisão?
					\item Verificar se o sistema pode ser tratado como isolado durante a colisão.
					\item Escrever os momentos lineares antes:
					\[
					\vec P_i=\sum m_i\vec v_{i}.
					\]
					\item Escrever os momentos lineares depois:
					\[
					\vec P_f=\sum m_i\vec v_{f}.
					\]
					\item Usar
					\[
					\textcolor{red}{\vec P_i=\vec P_f.}
					\]
					\item Se a colisão for elástica, acrescentar também a conservação da energia cinética (discutimos isso nas próximas slides).
				\end{enumerate}
			\end{column}
			
			\begin{column}{0.43\textwidth}
				\centering
				\begin{tikzpicture}[scale=0.72,>=stealth]
					\node[draw, rounded corners, fill=blue!5, align=center] (s1) at (0,3.2)
					{\scriptsize escolher\\ sistema};
					\node[draw, rounded corners, fill=blue!5, align=center] (s2) at (0,2.15)
					{\scriptsize verificar\\ isolamento};
					\node[draw, rounded corners, fill=blue!5, align=center] (s3) at (0,1.10)
					{\scriptsize escrever\\ estado inicial};
					\node[draw, rounded corners, fill=blue!5, align=center] (s4) at (0,0.05)
					{\scriptsize escrever\\ estado final};
					\node[draw, rounded corners, fill=red!5, draw=red!80!black, align=center] (s5) at (0,-1.10)
					{\scriptsize aplicar\\ \(\vec P_i=\vec P_f\)};
					\node[draw, rounded corners, fill=green!5, draw=green!60!black, align=center] (s6) at (0,-2.25)
					{\scriptsize se elástica:\\ \(K_i=K_f\)};
					
					\draw[->, thick] (s1) -- (s2);
					\draw[->, thick] (s2) -- (s3);
					\draw[->, thick] (s3) -- (s4);
					\draw[->, thick] (s4) -- (s5);
					\draw[->, thick] (s5) -- (s6);
				\end{tikzpicture}
			\end{column}
		\end{columns}
		
	\end{frame}
	
	
	% SLIDE 8
	\begin{frame}{\large Tipos de colisões: elástica, inelástica e perfeitamente inelástica}
		\tiny
		
		\textcolor{blue}{\textbf{Assumimos um sistema fechado e isolado durante a colisão.}}
		Assim, o momento linear total conserva-se:
		$
		\textcolor{red}{\vec P_i=\vec P_f.}
		$
		
		\vspace{0.15cm}
		
		\begin{center}
			\renewcommand{\arraystretch}{1.35}
			\begin{tabular}{c|c|c}
				\textbf{Tipo de colisão} & \textbf{Momento linear} & \textbf{Energia cinética} \\ \hline
				\textcolor{blue}{elástica} &
				\(\textcolor{red}{\vec P_i=\vec P_f}\) &
				\(\textcolor{red}{K_i=K_f}\) \\[2pt]
				
				\textcolor{blue}{inelástica} &
				\(\textcolor{red}{\vec P_i=\vec P_f}\) &
				\(\textcolor{red}{K_i\neq K_f}\), normalmente \(K_f<K_i\) \\[2pt]
				
				\textcolor{blue}{perfeitamente inelástica} &
				\(\textcolor{red}{\vec P_i=\vec P_f}\) &
				\(\textcolor{red}{K_f<K_i}\), corpos ficam juntos
			\end{tabular}
		\end{center}
		
		\vspace{0.15cm}
		
		\textcolor{red}{\textbf{Ideia-chave:}}
		na colisão elástica conservam-se \(\vec P\) e \(K\); na colisão inelástica conserva-se \(\vec P\), mas não \(K\); na perfeitamente inelástica, os corpos ficam unidos depois da colisão.
		
		\vspace{0.15cm}
		
		\begin{columns}[T]
			% Exemplo elástico
			\begin{column}{0.50\textwidth}
				\textcolor{blue}{\textbf{Exemplo simples: colisão elástica}}
				
				\vspace{0.1cm}
				
				Duas massas iguais; a segunda está inicialmente em repouso:
				$
				m_1=m_2=m,
				\qquad
				v_{1i}=v,
				\qquad
				v_{2i}=0.
				$
				\vspace{0.3cm}
				
				Considere uma colisão elástica frontal:
				$
				{v_{1f}=0,
					v_{2f}=v.}
				$ e $
				{K_i=K_f.}
				$
				
				\vspace{0.3cm}
				
				\centering
				\begin{tikzpicture}[scale=0.75,>=stealth]
					\node[font=\tiny] at (-0.3,1.55) {antes};
					\draw[thick] (-2.0,1.0) -- (2.0,1.0);
					\draw[fill=blue!20, draw=blue!80!black] (-1.05,1.0) circle (0.16);
					\draw[fill=green!20, draw=green!60!black] (0.75,1.0) circle (0.16);
					\draw[vetorv] (-1.05,1.0) -- (-0.25,1.0) node[midway, above] {\tiny \(v\)};
					\node[above] at (0.75,1.18) {\tiny repouso};
					
					\node[font=\tiny] at (0,0.25) {depois};
					\draw[thick] (-2.0,-0.25) -- (2.0,-0.25);
					\draw[fill=blue!20, draw=blue!80!black] (-0.45,-0.25) circle (0.16);
					\node[below] at (-0.45,-0.43) {\tiny para};
					\draw[fill=green!20, draw=green!60!black] (0.45,-0.25) circle (0.16);
					\draw[vetorv] (0.45,-0.25) -- (1.25,-0.25) node[midway, above] {\tiny \(v\)};
					
					\node[align=center, font=\tiny, red!80!black] at (0,-1.05)
					{o momento é totalmente transferido para a segunda massa \(m_2\)};
				\end{tikzpicture}
			\end{column}
			
			% Exemplo inelástico não perfeito
			\begin{column}{0.50\textwidth}
				\textcolor{blue}{\textbf{Exemplo simples: colisão inelástica}}
				
				\vspace{0.1cm}
				
				Uma bola colide com uma parede e ricocheteia, mas perde parte da energia cinética:
				${K_f<K_i}$. A bola não fica presa à parede:
				\[
				\textcolor{red}{\text{não é perfeitamente inelástica}.}
				\]
				
				O momento linear da bola muda:
				$
				\Delta \vec p=\vec p_f-\vec p_i.
				$
				
				\vspace{0.3cm}
				
				\centering
				\begin{tikzpicture}[scale=0.78,>=stealth]
					\draw[thick] (-2.0,-0.35) -- (1.8,-0.35);
					\draw[fill=gray!20, draw=black, pattern=north east lines] (0.95,-0.35) rectangle (1.25,1.25);
					\node[right] at (1.25,0.45) {\tiny parede};
					
					% antes
					\draw[fill=blue!25, draw=blue!80!black] (-1.15,0.25) circle (0.14);
					\draw[vetorv] (-1.15,0.25) -- (-0.35,0.25) node[midway, above] {\tiny \(\vec p_i\)};
					
					% depois
					\draw[fill=blue!25, draw=blue!80!black] (0.45,0.25) circle (0.14);
					\draw[vetorv] (0.45,0.25) -- (-0.15,0.25) node[midway, below] {\tiny \(\vec p_f\)};
					
					% perda de energia
					\node[align=center, font=\tiny, red!80!black] at (-0.25,-1.05)
					{ricocheteia, mas com\\ menor energia cinética};
				\end{tikzpicture}
				
				\vspace{0.05cm}
				
				\[
				\textcolor{red}{K_f<K_i, \qquad \text{corpos não ficam unidos}.}
				\]
			\end{column}
		\end{columns}
		
	\end{frame}
	
	
	% SLIDE 9
	\begin{frame}{Equações de conservação nas colisões}
		\scriptsize
		
		Considere dois corpos em movimento numa reta. Escolhemos o eixo \(x\) ao longo da reta da colisão.
		As velocidades são escritas como componentes:
		\[
		v_{1x,i},\quad v_{2x,i}
		\qquad \text{antes da colisão,}
		\qquad
		v_{1x,f},\quad v_{2x,f}
		\qquad \text{depois da colisão.}
		\]
		
		\textcolor{blue}{\textbf{Nota:}}
		o sinal de cada componente indica o sentido do movimento no eixo \(x\).
		Para uma colisão elástica, usamos duas leis:
		
		\begin{columns}[T]
			\begin{column}{0.58\textwidth}
				\[
				\textcolor{blue}{\textbf{Momento linear:}\,\,}
				\textcolor{red}{
					m_1v_{1x,i}+m_2v_{2x,i}
					=
					m_1v_{1x,f}+m_2v_{2x,f}.}
				\]
				
				\[
				\begin{aligned}
					\textcolor{blue}{\textbf{Energia cinética:}}\quad
					&\textcolor{red}{
						\frac12m_1v_{1x,i}^2+\frac12m_2v_{2x,i}^2
						=
						\frac12m_1v_{1x,f}^2+\frac12m_2v_{2x,f}^2.}
				\end{aligned}
				\]
				\vspace{0.1cm}
				{(A conservação da energia cinética aplica-se apenas às colisões elásticas.)}
				
			\end{column}
			
			\begin{column}{0.42\textwidth}
				\makebox[\linewidth][r]{%
					\begin{tikzpicture}[scale=0.75,>=stealth]
						\draw[eixo] (-2.25,2.05) -- (2.25,2.05) node[right] {\scriptsize \(x\)};
						
						\node[font=\scriptsize] at (0,2.55) {antes};
						\draw[thick] (-2.1,1.55) -- (2.1,1.55);
						
						\draw[fill=blue!25, draw=blue!80!black] (-1.25,1.55) circle (0.18);
						\node[below] at (-1.25,1.25) {\scriptsize \(m_1\)};
						\draw[vetorv] (-1.25,1.55) -- (-0.25,1.55)
						node[midway, above] {\scriptsize \(v_{1x,i}\)};
						
						\draw[fill=green!25, draw=green!60!black] (1.25,1.55) circle (0.14);
						\node[below] at (1.25,1.25) {\scriptsize \(m_2\)};
						\draw[vetorv] (1.25,1.55) -- (0.35,1.55)
						node[midway, above] {\scriptsize \(v_{2x,i}\)};
						
						\node[font=\scriptsize] at (0,0.65) {depois};
						\draw[thick] (-2.1,-0.25) -- (2.1,-0.25);
						
						\draw[fill=blue!25, draw=blue!80!black] (-0.45,-0.25) circle (0.18);
						\draw[vetorv] (-0.45,-0.25) -- (-1.25,-0.25)
						node[midway, below] {\scriptsize \(v_{1x,f}\)};
						
						\draw[fill=green!25, draw=green!60!black] (0.45,-0.25) circle (0.14);
						\draw[vetorv] (0.45,-0.25) -- (1.55,-0.25)
						node[midway, below] {\scriptsize \(v_{2x,f}\)};
						
						\node[align=center, font=\scriptsize, red!80!black] at (0,-1.1)
						{};
					\end{tikzpicture}%
				}
			\end{column}
		\end{columns}
		
		\vspace{0.25cm}
		
		\textcolor{blue}{\textbf{Como usar na prática:}}
		com os dados do problema, substituímos os valores conhecidos nas duas equações acima.
		Se houver duas incógnitas, o sistema permite determinar as massas ou velocidades desconhecidas.
		
	\end{frame}
	
	
	% SLIDE 10
	\begin{frame}{Exemplos: colisão elástica e perfeitamente inelástica}
		\tiny
		
		\begin{columns}[T]
			% COLUNA ESQUERDA: COLISÃO ELÁSTICA
			\begin{column}{0.50\textwidth}
				\textcolor{blue}{\textbf{Colisão elástica frontal com alvo em repouso}}
				
				\vspace{0.1cm}
				
				Projétil \(m_1\) com velocidade inicial \(v_{1i}\) colide elasticamente com um alvo \(m_2\) inicialmente em repouso:
				$
				\textcolor{red}{v_{2i}=0.}
				$
				
				\vspace{0.1cm}
				
				Da conservação do momento linear e da energia cinética:
				\[
				m_1v_{1i}=m_1v_{1f}+m_2v_{2f},
				\qquad
				\frac12m_1v_{1i}^2=
				\frac12m_1v_{1f}^2+\frac12m_2v_{2f}^2.
				\]
				
				Da primeira equação:
				$
				\textcolor{red}{m_1(v_{1i}-v_{1f})=m_2v_{2f}.}
				$
				\vspace{0.1cm}
				
				Da segunda:
				$
				\textcolor{red}{m_1(v_{1i}^2-v_{1f}^2)=m_2v_{2f}^2.}
				$
				
				\vspace{0.1cm}
				
				Como \(v_{1i}^2-v_{1f}^2=(v_{1i}-v_{1f})(v_{1i}+v_{1f})\), dividindo as duas expressões:
				\[
				\textcolor{red}{v_{2f}=v_{1i}+v_{1f}.}
				\]
				
				Substituindo:
				$
				m_1(v_{1i}-v_{1f})=m_2(v_{1i}+v_{1f})
				$
				\[
				\textcolor{red}{
					v_{1f}=
					\frac{m_1-m_2}{m_1+m_2}\,v_{1i}},
				\qquad
				\textcolor{red}{
					v_{2f}=
					\frac{2m_1}{m_1+m_2}\,v_{1i}.}
				\]
				
				Se \(m_1=m_2\), então
				$
				\textcolor{red}{v_{1f}=0,\quad v_{2f}=v_{1i}.}
				$
				
				\centering
				\begin{tikzpicture}[scale=0.62,>=stealth]
					\draw[thick] (-2.1,1.5) -- (2.1,1.5);
					\node[font=\tiny] at (-1.55,2.15) {antes};
					
					\draw[fill=blue!25, draw=blue!80!black] (-1.2,1.5) circle (0.17);
					\draw[vetorv] (-1.05,1.5) -- (-0.25,1.5) node[midway, above] {\tiny \(v_{1i}\)};
					
					\draw[fill=green!25, draw=green!60!black] (0.9,1.5) circle (0.17);
					\node[above] at (0.9,1.72) {\tiny \(v_{2i}=0\)};
					
					\draw[thick] (-2.1,0.0) -- (2.1,0.0);
					\node[font=\tiny] at (-1.05,0.65) {depois, se \(m_1=m_2\)};
					
					\draw[fill=blue!25, draw=blue!80!black] (-0.35,0) circle (0.17);
					\node[below] at (-0.35,-0.22) {\tiny para};
					
					\draw[fill=green!25, draw=green!60!black] (0.75,0) circle (0.17);
					\draw[vetorv] (0.9,0) -- (1.75,0) node[midway, above] {\tiny \(v_{1i}\)};
					
					\node[align=center, font=\tiny, red!80!black] at (0,-0.9)
					{};
				\end{tikzpicture}
			\end{column}
			
			% COLUNA DIREITA: COLISÃO PERFEITAMENTE INELÁSTICA
			\begin{column}{0.50\textwidth}
				\textcolor{blue}{\textbf{Colisão perfeitamente inelástica: carruagens}}
				
				\vspace{0.1cm}
				
				Uma carruagem aproxima-se de outra em repouso com
				$
				v=0{,}5\,\text{m/s}.
				$
				Depois da colisão, as carruagens acoplam-se e movem-se juntas com a mesma velocidade \(v'\).
				
				\vspace{0.1cm}
				
				Se as carruagens têm a mesma massa \(m\):
				\[
				mv+m\cdot 0=mv'+mv'
				\]
				\[
				mv=2mv'
				\qquad \Rightarrow \qquad
				\textcolor{red}{v'=\frac{v}{2}=0{,}25\,\text{m/s}.}
				\]
				
				A energia cinética não se conserva:
				\[
				K_i=\frac12mv^2,
				\qquad
				K_f=\frac12(2m)\left(\frac v2\right)^2=\frac14mv^2.
				\]
				Logo,
				$
				\textcolor{red}{K_f<K_i.}
				$
				
				\centering
				\begin{tikzpicture}[scale=0.68,>=stealth]
					\draw[thick] (-2.2,2.0) -- (2.2,2.0);
					\node[font=\tiny] at (-1.55,2.75) {antes};
					
					\draw[fill=blue!10, draw=blue!80!black] (-1.85,2.0) rectangle (-0.85,2.55);
					\draw[fill=blue!10, draw=blue!80!black] (0.65,2.0) rectangle (1.65,2.55);
					\draw[vetorv] (-0.85,2.28) -- (-0.1,2.28) node[midway, above] {\tiny \(v\)};
					\node[above] at (1.15,2.55) {\tiny repouso};
					
					\draw[thick] (-2.2,0.2) -- (2.2,0.2);
					\node[font=\tiny] at (0.55,1.0) {depois};
					
					\draw[fill=green!10, draw=green!60!black] (-0.55,0.2) rectangle (0.45,0.75);
					\draw[fill=green!10, draw=green!60!black] (0.45,0.2) rectangle (1.45,0.75);
					\draw[vetorv] (1.45,0.48) -- (2.05,0.48) node[midway, above] {\tiny \(v'\)};
					
					\node[align=center, font=\tiny, red!80!black] at (0,-0.75)
					{corpos acoplados:\\ colisão perfeitamente inelástica};
				\end{tikzpicture}
			\end{column}
		\end{columns}
		
	\end{frame}
	
	
	% SLIDE 11
	\begin{frame}{Colisões em duas dimensões}
		\scriptsize
		
		Se os corpos não se movem todos ao longo da mesma reta, a colisão é bidimensional. A conservação do momento continua válida para um sistema isolado:
		\[
		\textcolor{red}{\vec p_{1i}+\vec p_{2i}=\vec p_{1f}+\vec p_{2f}.}
		\]
		
		\begin{columns}[T]
			\begin{column}{0.58\textwidth}
				Na forma de componentes:
				\[
				\textcolor{red}{P_{x,i}=P_{x,f}},
				\qquad
				\textcolor{red}{P_{y,i}=P_{y,f}.}
				\]
				
				Para um alvo inicialmente em repouso e projétil inicialmente no eixo \(x\):
				\[
				x: \,\, m_1v_{1i}
				=
				m_1v_{1f}\cos\theta_1+
				m_2v_{2f}\cos\theta_2,
				\]
				\[
				y: \,\, 0
				=
				m_1v_{1f}\sin\theta_1-
				m_2v_{2f}\sin\theta_2.
				\]
				
				Se a colisão também for elástica:
				\[
				\textcolor{red}{
					\frac12m_1v_{1i}^2
					=
					\frac12m_1v_{1f}^2+
					\frac12m_2v_{2f}^2.}
				\]
			\end{column}
			
			\begin{column}{0.42\textwidth}
				\centering
				\begin{tikzpicture}[scale=0.8,>=stealth]
					\draw[eixo] (-1.8,0) -- (2.6,0) node[right] {$x$};
					\draw[eixo] (0,-1.7) -- (0,2.0) node[above] {$y$};
					
					\draw[fill=blue!25, draw=blue!80!black] (-1.4,0) circle (0.14);
					\draw[vetorv] (-1.25,0) -- (-0.25,0) node[midway, above] {\scriptsize \(v_{1i}\)};
					
					\draw[fill=green!25, draw=green!60!black] (0,0) circle (0.14);
					\node[below left] at (0,0) {\scriptsize \(m_2\)};
					
					\draw[fill=blue!25, draw=blue!80!black] (1.1,1.05) circle (0.14);
					\draw[vetorv] (0.15,0.1) -- (1.1,1.05) node[midway, above left] {\scriptsize \(v_{1f}\)};
					\draw[dashed] (0.25,0) -- (1.3,0);
					\node at (0.75,0.22) {\scriptsize \(\theta_1\)};
					
					\draw[fill=green!25, draw=green!60!black] (1.4,-0.95) circle (0.14);
					\draw[vetorv] (0.15,-0.1) -- (1.4,-0.95) node[midway, below left] {\scriptsize \(v_{2f}\)};
					\node at (0.85,-0.23) {\scriptsize \(\theta_2\)};
					
					\node[align=center, font=\scriptsize, red!80!black] at (0,-2.15)
					{duas equações de momento:\\ uma em \(x\), outra em \(y\)};
				\end{tikzpicture}
			\end{column}
		\end{columns}
		
	\end{frame}
	
	
	% SLIDE 12
	\begin{frame}{Exemplo: colisão de uma bola com uma parede}
		\scriptsize
		
		Considere uma bola que colide com uma parede. Durante o intervalo de contacto
		\(\Delta t\), a parede exerce uma força sobre a bola, alterando a componente do
		momento linear perpendicular à parede assim que essa componente muda so o sentido mas nao a magnitude.
		
		\vspace{0.2cm}
		
		Neste exemplo idealizado, considera-se que a componente perpendicular à parede
		troca de sinal, enquanto a componente paralela permanece igual.
		
		\vspace{0.2cm}
		
		Para a bola:
		$
		\textcolor{red}{p'_x=-p_x, \qquad p'_y=p_y.}
		$
		
		\vspace{0.2cm}
		
		Portanto,
		$
		\textcolor{red}{\Delta p_x=p'_x-p_x=-2p_x=-2mv_x,}
		\qquad
		\textcolor{red}{\Delta p_y=0.}
		$
		
		\vspace{0.2cm}
		
		\begin{columns}[T]
			\begin{column}{0.72\textwidth}
				
				A variação do momento linear da bola é causada pela força exercida pela
				parede durante o intervalo de contacto. Para o módulo da força média exercida pela parede sobre a bola:
				$
				\textcolor{red}{F_{\text{méd}}=\frac{|\Delta \vec p|}{\Delta t}.}
				$
				
				\vspace{0.35cm}
				
				Como apenas a componente \(x\) varia:
				$
				\textcolor{red}{|\Delta \vec p|=2|p_x|=2m|v_x|.}
				$
				
				\vspace{0.35cm}
				
				No caso de incidência normal:
				$
				\textcolor{red}{F_{\text{méd}}=\frac{2mv}{\Delta t}.}
				$
				
			\end{column}
			
			\begin{column}{0.28\textwidth}
				\centering
				
				\begin{tikzpicture}[scale=0.7,>=stealth]
					% parede
					\draw[thick, pattern=north east lines] (1.55,-1.5) rectangle (1.8,1.6);
					\node[right] at (1.8,1.45) {\scriptsize parede};
					
					% linha perpendicular no ponto de contacto
					\draw[dashed] (1.39,-1.2) -- (1.39,1.2);
					
					% bola antes da colisão
					\draw[fill=blue!25, draw=blue!80!black] (0.20,0.75) circle (0.14);
					\draw[vetorv] (0.34,0.70) -- (1.22,0.30)
					node[midway, above] {\scriptsize \(\vec v_i\)};
					
					% bola no instante do choque: encostada à parede
					\draw[fill=blue!25, draw=blue!80!black] (1.39,0.20) circle (0.16);
					\node[above left] at (1.32,0.38) {};
					
					% bola após ricochetear
					\draw[fill=blue!25, draw=blue!80!black] (0.45,-0.55) circle (0.14);
					\draw[vetorv] (1.23,0.08) -- (0.58,-0.48)
					node[midway, below] {\scriptsize \(\vec v_f\)};
					
					\node[align=center, font=\scriptsize, red!80!black] at (0,-1.95)
					{\(p_x\) troca de sinal;\\ \(p_y\) permanece};
				\end{tikzpicture}
			\end{column}
		\end{columns}
		
		\vspace{0.3cm}
		
		\textcolor{blue}{\textbf{Aplicação numérica:}}
		$
		m=60\,\text{g}=0{,}06\,\text{kg},
		\quad
		v=72\,\text{km/h}=20\,\text{m/s},
		\quad
		\Delta t=0{,}01\,\text{s},
		$
		
		\[
		F_{\text{méd}}=
		\frac{2(0{,}06)(20)}{0{,}01}
		=240\,\text{N}.
		\]
		
	\end{frame}
	
	
	% SLIDE 13
	\begin{frame}{Resumo para estudo}
		\scriptsize
		
		\begin{columns}[T]
			\begin{column}{0.58\textwidth}
				\textcolor{blue}{\textbf{Fórmulas centrais}}
				
				\[
				\vec r_{\text{cm}}=
				\frac{m_1\vec r_1+\cdots+m_n\vec r_n}{M},
				\qquad
				M=m_1+\cdots+m_n.
				\]
				
				\[
				M\vec a_{\text{cm}}=\vec F_{\text{ext}}.
				\]
				
				\[
				\vec p=m\vec v,
				\qquad
				\vec P=\vec p_1+\cdots+\vec p_n=M\vec v_{\text{cm}}.
				\]
				
				\[
				\vec F_{\text{ext}}=0
				\qquad \Rightarrow \qquad
				\vec P_i=\vec P_f.
				\]
				
				\[
				\Delta \vec p=\vec p_f-\vec p_i,
				\qquad
				F=\frac{\Delta p}{\Delta t}.
				\]
			\end{column}
			
			\begin{column}{0.42\textwidth}
				\textcolor{blue}{\textbf{Perguntas de controlo}}
				
				\begin{enumerate}
					\item Qual é o sistema escolhido?
					\item A força externa resultante é nula?
					\item O problema é unidimensional ou bidimensional?
					\item Os corpos ficam juntos depois da colisão?
					\item A colisão é elástica?
					\item Devo usar só \(\vec P_i=\vec P_f\), ou também \(K_i=K_f\)?
				\end{enumerate}
				
				\vspace{0.15cm}
				
				\textcolor{red}{\textbf{Mensagem final:}}
				
				O centro de massa descreve o movimento global do sistema; o momento linear é a grandeza conservada quando não há força externa resultante.
			\end{column}
		\end{columns}
		
		\vspace{0.3cm}
		
		{
			\textcolor{red}{\textbf{Atenção:}}
			se existe força externa resultante, por exemplo atrito, então
			$
			\textcolor{red}{
				\frac{d\vec P}{dt}=\vec F_{\text{ext}}\neq 0.}
		$
			Nesse caso, em geral, \textcolor{red}{não podemos usar} \(\vec P_i=\vec P_f\).
		}
		
	\end{frame}
	
\end{document}
